\section{Detailed proofs}\label{app:proofs}

This appendix proves the claims in the main text with one concrete choice of the scale constants.
No optimization of numerical constants is intended.

\subsection{Notation and elementary bounds}

Let $\varepsilon_1,\ldots,\varepsilon_n$ be independent Rademacher signs and
$B_n=\sum_i\varepsilon_i$. Define
\[
 k_\star=\left\lceil\frac{5n}{8}\right\rceil,
 \qquad
 s_\star=2k_\star-n,
 \qquad
 q_\star=\Pp(B_n\ge s_\star).
\]
Then
\begin{equation}
 \frac14\le\frac{s_\star}{n}\le\frac14+\frac2n.
 \label{eq:rho-range}
\end{equation}
Hoeffding's inequality gives
\begin{equation}
 q_\star\le \exp\!\left(-\frac{s_\star^2}{2n}\right)
 \le e^{-n/32}.
 \label{eq:qstar-hoeffding}
\end{equation}

For $r\le n/1024$, put
\[
 q(r)=\Pp(B_n\ge s_\star-2r).
\]
The event on the right contains the upper tail at $k_\star$ and at most $r$ additional binomial
levels. Moving down one level multiplies the point mass by
\[
 \frac{k}{n-k+1}
 \le \frac{5n/8+1}{3n/8-r}
 <2
\]
for all sufficiently large $n$. Since $q_\star$ is at least the mass at $k_\star$, we obtain
\begin{equation}
 \frac{q(r)}{q_\star}
 \le(r+1)2^r\le4^r.
 \label{eq:tail-ratio}
\end{equation}
The last inequality holds for every integer $r\ge1$.

\subsection{The symmetrized maximum}

\begin{proof}[Proof of Lemma~\ref{lem:H}]
Write $M_a(x)=\max_ja_jx_j$. Then
\[
 H_a(x)=M_a(x)-M_a(-x).
\]
This immediately gives $H_a(-x)=-H_a(x)$. Since each $a_jx_j$ lies in
$[-\|a\|_\infty,\|a\|_\infty]$, the range bound follows. Finally,
\[
 |M_a(x)-M_{a'}(x)|\le\max_j|(a_j-a'_j)x_j|
 \le\|a-a'\|_\infty.
\]
Applying this inequality at $x$ and $-x$ and using the triangle inequality proves the final claim.
\end{proof}

\subsection{The learner is stable and bounded}\label{app:stable}

We verify the claims around \eqref{eq:predictor}. Suppose $S$ and $S'$ differ at observation $i$.
For every ramp coordinate,
\[
 |S_{kj}-S'_{kj}|=|X_{i,kj}-X'_{i,kj}|\le2.
\]
The scalar map in \eqref{eq:ramp} is $\gamma/2$-Lipschitz. Therefore
\begin{equation}
 \|a(S)-a(S')\|_\infty\le\gamma.
 \label{eq:a-stability}
\end{equation}
Lemma~\ref{lem:H} implies that the $H/4$ term changes by at most $\gamma/2$ at every test point.

Only the old and new memory cells can change. At either cell, $b_j$ takes values in
$\{-c_{\rm mem},0,c_{\rm mem}\}$, so
\begin{equation}
 |b_j(S)-b_j(S')|\le2c_{\rm mem}\le\gamma/2.
 \label{eq:b-stability}
\end{equation}
For a fixed test point, only its own cell is read. Adding \eqref{eq:a-stability} and
\eqref{eq:b-stability} proves $\gamma$-uniform stability.

If the scale set is nonempty, its largest element satisfies $p_K\le L/\gamma$. Thus
\begin{equation}
 \|a(S)\|_\infty\le\gamma r_K=\frac{\gamma p_K}{16}\le\frac L{16}.
 \label{eq:a-range}
\end{equation}
The same bound is trivial if there are no ramps. Lemma~\ref{lem:H} gives
$|H_a(X)|/4\le L/32$. Since $|b_J\Sigma|\le L/8$ and $|LU/4|=L/4$,
\[
 3L/32\le h_S(Z)\le29L/32.
\]
This proves the range claim. With regression label zero, absolute loss is exactly $h_S(Z)$.

For a fresh input, $X\overset{d}{=}-X$, and Lemma~\ref{lem:H} gives
$\E[H_{a(S)}(X)\mid S]=0$. The fresh signs $\Sigma$ and $U$ are independent of $S$ and centered.
This proves \eqref{eq:population-half} and \eqref{eq:gap-decomp}.

\subsection{Probability of a clean extreme}\label{app:extreme}

We prove the probability part of Lemma~\ref{lem:extreme}. For a scale $p_k=16r_k\le n/64$,
equation \eqref{eq:qstar-hoeffding} gives
\begin{equation}
 q_\star\le e^{-2p_k}.
 \label{eq:qstar-small}
\end{equation}
By the definition of $d_k$,
\begin{align}
 d_kq_\star
 &\ge 8e^{-p_k/2}-q_\star
 \ge7e^{-p_k/2},
 &d_kq_\star&\le8e^{-p_k/2}.
 \label{eq:lambda-bracket}
\end{align}
In particular, $d_k\ge2$.

The clean event forbids every nonexceptional coordinate whose score is at least
$t_k=s_\star-2r_k$. Equations \eqref{eq:tail-ratio} and \eqref{eq:lambda-bracket} show
\begin{equation}
 d_k q(r_k)
 \le8e^{-p_k/2}4^{r_k}
 \le8e^{-3p_k/8},
 \label{eq:active-expectation}
\end{equation}
because $4^{r_k}=e^{(\log4)p_k/16}<e^{p_k/8}$. The same inequality holds at every larger scale.
Since $p_{k+j}=4^jp_k$,
\begin{equation}
 \sum_{\ell\ge k}d_\ell q(r_\ell)
 \le8\sum_{j\ge0}\exp\!\left(-\frac38 4^jp_k\right)
 <0.021
 \label{eq:forbidden-sum}
\end{equation}
for $p_k\ge16$.

Choose which coordinate in group $k$ is the exceptional one. Independence of all Rademacher
coordinates gives
\begin{align}
 \Pp(E_k)
 &=d_kq_\star
   (1-q(r_k))^{d_k-1}
   \prod_{\ell>k}(1-q(r_\ell))^{d_\ell}\\
 &\ge d_kq_\star
 \left(1-\sum_{\ell\ge k}d_\ell q(r_\ell)\right)\\
 &\ge7e^{-p_k/2}(1-0.021)
 \ge6e^{-p_k/2}.
 \label{eq:Ek-proof}
\end{align}
The second line is the union bound applied to all forbidden coordinates. This proves
\eqref{eq:extreme-prob}.

\subsection{Gap on the clean extreme event}\label{app:gap}

On $E_k$, let $(k,j_\star)$ denote the exceptional coordinate and put $a=\gamma r_k$. Its score
is at least $s_\star=t_k+2r_k$, so its ramp reaches the cap $a$. Every other coordinate in group
$k$ or a larger group has amplitude zero. Every smaller scale has cap at most $a/4$.

Write $x_i^\star=X_{i,kj_\star}$ and let $n_+$ and $n_-$ count its positive and negative signs.
There is at least one zero-amplitude coordinate because $d_k\ge2$. If $x_i^\star=1$, the first
maximum in $H$ is $a$ and the second is at most $a/4$, so
\begin{equation}
 H_{a(S)}(X_i)\ge3a/4.
 \label{eq:H-plus}
\end{equation}
If $x_i^\star=-1$, the second maximum is $a$ and the first is at least zero, so
\begin{equation}
 H_{a(S)}(X_i)\ge-a.
 \label{eq:H-minus}
\end{equation}
Let $\rho=(n_+-n_-)/n$. Equations \eqref{eq:H-plus} and \eqref{eq:H-minus} imply
\begin{align}
 \frac1n\sum_iH_{a(S)}(X_i)
 &\ge \frac{3a}{4}\frac{1+\rho}{2}-a\frac{1-\rho}{2}
 =a\left(\frac78\rho-\frac18\right).
 \label{eq:H-average-rho}
\end{align}
On $E_k$, $\rho\ge s_\star/n\ge1/4$. The right side of \eqref{eq:H-average-rho} is increasing in
$\rho$, so it is at least $3a/32$. Substituting $a=\gamma r_k=\gamma p_k/16$ proves
\eqref{eq:extreme-gap}.

\subsection{Rademacher anti-concentration}\label{app:anti}

We give an elementary block proof of Lemma~\ref{lem:rad}. If $V_m$ is a sum of $m$ independent
Rademacher signs, then
\[
 \E V_m^2=m,
 \qquad
 \E V_m^4=3m^2-2m\le3m^2.
\]
Paley--Zygmund applied to $V_m^2$ therefore gives
\begin{equation}
 \Pp\!\left(|V_m|\ge\sqrt{m/2}\right)\ge\frac1{12}.
 \label{eq:block-paley}
\end{equation}
By symmetry, either one-sided event has probability at least $1/24$.

Partition the first $qm$ signs of $T_n$ into $q$ blocks of size $m=\lfloor n/q\rfloor$, leaving
at most $q-1$ signs. If every block sum is at least $\sqrt{m/2}$ and the leftover sum is
nonnegative, then, whenever $q\le n/2$,
\begin{equation}
 T_n\ge q\sqrt{m/2}\ge\frac12\sqrt{nq}.
 \label{eq:block-size}
\end{equation}
The event has probability at least
\begin{equation}
 \frac12\,24^{-q}.
 \label{eq:block-prob}
\end{equation}
We now track constants. First suppose $1\le p<16$. The largest atom of $T_n$ is at most
$n^{-1/2}$ by the central binomial coefficient bound. Positive attainable values are spaced by
two, so there are at most $\sqrt n/32+1$ of them below $\sqrt n/16$. Symmetry gives
\begin{equation}
 \Pp(T_n\ge\sqrt n/16)
 \ge\frac12-\frac{1}{2\sqrt n}-\left(\frac1{32}+\frac1{\sqrt n}\right)
 =\frac{15}{32}-\frac{3}{2\sqrt n}\ge0.45
 \label{eq:small-rad}
\end{equation}
for all sufficiently large $n$. Since $p<16$, the threshold is at least
$\sqrt{np}/64$.

For $16\le p<32\log24$, use one block. Equations~\eqref{eq:block-paley} and symmetry give
\[
 \Pp(T_n\ge\sqrt{n/2})\ge1/24\ge(1/8)e^{-p/8},
\]
and $\sqrt{n/2}\ge\sqrt{np}/64$. Finally, suppose $32\log24\le p\le n$ and take
\[
 q=\left\lfloor\frac{p}{16\log24}\right\rfloor.
\]
Then $p/(32\log24)\le q\le p/(16\log24)$ and $q\le n/2$. Equations
\eqref{eq:block-size} and \eqref{eq:block-prob} yield
\[
 T_n\ge\frac{1}{2\sqrt{32\log24}}\sqrt{np}
 \ge\frac1{64}\sqrt{np}
\]
with probability at least
\[
 \frac12 24^{-q}\ge\frac12e^{-p/16}
 \ge\frac18e^{-p/8}.
\]
The three ranges prove Lemma~\ref{lem:rad} with no asymptotic normal approximation.

\subsection{Small scales, saturation, and simultaneity}\label{app:cases}

We fill in the cases abbreviated in the main proof. First, the memory indices are distinct with
probability at least $15/16$ by \eqref{eq:distinct}. On this event,
\begin{equation}
 \frac1n\sum_i b_{J_i}(S)\Sigma_i=c_{\rm mem}.
 \label{eq:memory-gap}
\end{equation}

Suppose $1\le p<16$. If $\gamma\le L/2$, then $c_{\rm mem}=\gamma/4$ and
\[
 c_{\rm mem}\ge\frac1{64}\min\{L,\gamma p\}.
\]
If $\gamma>L/2$, then $c_{\rm mem}=L/8$ and the same inequality is immediate. The probability
that any ramp is active is at most the sum in \eqref{eq:forbidden-sum}, hence below $0.021$.
By the compact-interval part of Lemma~\ref{lem:rad}, the sampling event can be chosen to have
probability at least $0.45$ while still giving a universal multiple of $L\sqrt{p/n}$.
The three events are independent, and
$0.979\cdot(15/16)\cdot0.45>e^{-1}$. Their intersection therefore exceeds $e^{-p}$.
This proves \eqref{eq:main-tail} for $p<16$.

Now suppose $p\ge16$ but $L/\gamma<16$. Then $\gamma>L/16$ and
\begin{equation}
 c_{\rm mem}=\min\{\gamma/4,L/8\}\ge L/64.
 \label{eq:memory-saturated}
\end{equation}
Thus memory alone supplies a constant fraction of the saturated stability term. Intersecting
\eqref{eq:memory-gap} with Lemma~\ref{lem:rad} proves the result.

It remains to justify \eqref{eq:chosen-scale}. Let
$P=\min\{n/64,L/\gamma\}$. If $P\ge16$, the largest geometric scale below $P$ is greater than
$P/4$. For a query $p\le c_0n\le n/64$, the largest available scale not exceeding $p$, or the terminal
scale if $p>P$, therefore obeys
\[
 p_k\ge\frac14\min\{p,L/\gamma\}.
\]
The construction of all groups uses only $(n,L,\gamma)$. It is completed before $p$ is chosen.
Thus the proof supplies the probability statement for all $p$ using the same distribution and
algorithm. This establishes the simultaneity claimed in Theorem~\ref{thm:main}.

Finally, on the intersection used in the proof,
\begin{align}
 G_S
 &\ge \frac14\frac{3\gamma p_k}{512}
       +c_{\rm mem}
       +\frac{cL}{4}\sqrt{p/n}\\
 &\ge c'\left(\min\{L,\gamma p\}+L\sqrt{p/n}\right)\\
 &\ge c'\min\!\left\{L,\gamma p+L\sqrt{p/n}\right\}.
\end{align}
This also makes clear that every term is one-sided and no cancellation is used.

\subsection{Moment consequence}

Let
\[
 u_p=c_1\min\!\left\{L,\gamma p+L\sqrt{p/n}\right\}.
\]
Theorem~\ref{thm:main} gives $\Pp(G_S\ge u_p)\ge e^{-p}$. Hence
\[
 \|G_S\|_p^p\ge u_p^p e^{-p},
 \qquad
 \|G_S\|_p\ge e^{-1}u_p.
\]
This proves Corollary~\ref{cor:moments}. The upper direction of \eqref{eq:minimax} is
\eqref{eq:upper} combined with $|G_S|\le L$. The lower direction is the corollary.

\section{Exact computational audits}\label{app:audit}

The construction contains only binomial probabilities and maximum inequalities. We nevertheless
performed three audits.

\paragraph{Exact clean-event probabilities.}
For each $n\in\{4096,8192,16384,32768,65536\}$ and every geometric scale below $n/64$, we
evaluated $q_\star$, each lowered tail $q(r_k)$, the floor in $d_k$, and the exact product in
\eqref{eq:Ek-proof} in log space. Table~\ref{tab:audit} shows representative results. The final
column divides the rigorous clean-event lower bound by $e^{-p_k/2}$. The analytic proof uses the
weaker constant $6$.

\begin{table}[h]
\centering
\caption{Exact audit of the multiscale event.}
\label{tab:audit}
\begin{tabular}{rrrrr}
\toprule
$n$ & $p_k$ & $-\log q_\star$ & $\log(q(r_k)/q_\star)/p_k$ &
$\Pp(E_k)e^{p_k/2}$ lower bound\\
\midrule
4096  & 16   & 132.81  & 0.0320 & 7.964\\
4096  & 64   & 132.81  & 0.0319 & 8.000\\
16384 & 256  & 521.60  & 0.0318 & 8.000\\
65536 & 1024 & 2074.71 & 0.0318 & 8.000\\
\bottomrule
\end{tabular}
\end{table}

\paragraph{Exhaustive stability audit.}
For $n=d=3$, we enumerated every Rademacher dataset, every one-point replacement, and every test
input. We separately enumerated every memory dataset with three cells. With $\gamma=0.2$, the
largest ramp loss change was $0.1$, the largest memory loss change was $0.1$, and their sum was
exactly $0.2$. The maximum coordinate-amplitude change was exactly $\gamma$. This checks the
sharp cases of \eqref{eq:a-stability} and \eqref{eq:b-stability}.

The audit scripts use exact enumeration for stability and double-precision log binomial masses for
the probability calculation. The displayed margins are several orders of magnitude larger than
floating-point error.
